% Self-contained proof of the exact fixed-alpha worst-list bound.
\documentclass[11pt]{article}

\usepackage[T1]{fontenc}
\usepackage[utf8]{inputenc}
\usepackage{lmodern}
\usepackage{amsmath,amssymb,amsthm,mathtools}
\usepackage[margin=1in]{geometry}
\usepackage{microtype}
\usepackage[colorlinks=true,linkcolor=blue,citecolor=blue,urlcolor=blue]{hyperref}

\newtheorem{theorem}{Theorem}[section]
\newtheorem{proposition}[theorem]{Proposition}
\newtheorem{corollary}[theorem]{Corollary}
\theoremstyle{remark}
\newtheorem{remark}[theorem]{Remark}

\newcommand{\OPT}{\operatorname{OPT}}
\newcommand{\WLPT}{\operatorname{WLPT}}

\title{The Exact Worst-List Bound for
\texorpdfstring{$P\,\|\,\sum_j p_j C_j(\alpha)$}
{P || sum p_j Cj(alpha)}
for \texorpdfstring{$0<\alpha\leq\tfrac12$}{0 < alpha <= 1/2}}
\author{}
\date{}

\begin{document}
\maketitle

\begin{abstract}
For jobs of weight $w_j=p_j$ on identical parallel machines, let
$C_j(\alpha)=S_j+\alpha p_j$, where $0<\alpha\leq1/2$.  This note
proves, self-containedly, that the exact worst-list approximation ratio is
\[
 1+\frac{1}{2\alpha H(\alpha)},\qquad
 H(\alpha)=\min_{k\geq2}\left\{
 \frac{k+1}{2k}+
 \sqrt{\left(\frac{k+1}{2k}\right)^2+
 k+\frac1k+\frac1\alpha}\right\}.
\]
It is strictly smaller than $1+1/(6\alpha)$ for
$0<\alpha<1/2$ and equals the tight value $4/3$ for this worst-list
rule at $\alpha=1/2$.  A finite matching instance is given for every
$\alpha$.
\end{abstract}

\section{Model, convention, and upper-bound theorem}

All jobs have processing time $p_j>0$, weight $w_j=p_j$, release time
zero, and are processed nonpreemptively on identical machines.  In this
note, the symbol \emph{WLPT} is used nonstandardly for the worst
admissible list realization: jobs are taken from an arbitrary list and
each next job is assigned to a least-loaded machine.  Thus
\[
 \WLPT_\alpha(I):=\max_W\sum_jp_jC_j^W(\alpha),
 \qquad
 \OPT_\alpha(I):=\min_S\sum_jp_jC_j^S(\alpha),
\]
where the maximum includes minimum-load tie choices.  This convention is
natural here because all Smith ratios $w_j/p_j$ are equal, so the
unrestricted priority-tie family of the WSPT/Smith rule is exactly the
family of all lists.  The upper bound below applies to every individual
list and every tie choice.

\begin{theorem}\label{thm:exact-wlpt-alpha}
Consider finitely many jobs with processing times \(p_j>0\), zero release
times, and nonpreemptive processing on identical parallel machines.  Write
\[
 C_j(\alpha):=S_j+\alpha p_j,\qquad 0<\alpha\leq\frac12,
\]
where \(S_j\) is the start time.  A \emph{list schedule} processes an
arbitrary ordered list of the jobs and assigns each successive job to a
machine of minimum current load; any minimum-load machine may be chosen
when there is a tie.  For $k\in\mathbb N$, put
\[
 b_k:=\frac{k+1}{2k},\qquad
 h_k:=b_k+\sqrt{b_k^2+k+\frac1k+\frac1\alpha},
 \qquad
 \rho_k:=1+\frac{1}{2\alpha h_k},
\]
and let
\[
 H(\alpha):=\min_{k\geq2}h_k,
 \qquad
 \rho=\rho(\alpha):=\max_{k\geq1}\rho_k
 =1+\frac{1}{2\alpha H(\alpha)}.
\]
Then every such list schedule $W$ satisfies
\[
 \sum_jp_jC_j^W(\alpha)
 \leq \rho(\alpha)\min_S\sum_jp_jC_j^S(\alpha).
\]
The maximum in the definition of $\rho$ is attained.
\end{theorem}

\begin{proof}
All schedules below are taken to be left-justified.  This is automatic for
a list schedule, and an optimum can be left-shifted without increasing its
objective.  For a left-justified assignment $S$, write
\[
 Q_J:=\sum_jp_j^2,
 \qquad
 C(S):=\sum_i\sum_{\{j,\ell\}\subseteq S_i}p_jp_\ell.
\]
The order of the jobs on a fixed machine does not affect the objective,
because every unordered pair on that machine contributes its product
exactly once.  Consequently,
\begin{equation}\label{eq:objective-pair-form}
 \gamma_{\alpha}(S):=\sum_jp_jC_j^S(\alpha)=\alpha Q_J+C(S).
\end{equation}
If $L_i(S)$ is the load of machine $i$, then
$C(S)=\tfrac12(\sum_iL_i(S)^2-Q_J)$, and hence also
\begin{equation}\label{eq:objective-load-form}
 \gamma_{\alpha}(S)
 =\frac12\sum_iL_i(S)^2+
   \left(\alpha-\frac12\right)Q_J.
\end{equation}

Since $h_k\to\infty$, one has $\rho_k\to1$; hence the maximum of
the sequence $\{\rho_k\}$ is attained.  Choose an index $n$ attaining
it, and define
\begin{equation}\label{eq:main-parameters}
 h:=h_n,
 \quad D:=(\rho-1)\alpha=\frac1{2h},
 \quad a:=\frac\rho2,
 \quad A:=\frac12+D,
 \quad \beta:=a-A=(\rho-1)\left(\frac12-\alpha\right)\geq0,
\end{equation}
as well as
\begin{equation}\label{eq:C0-M-K0}
 C_0:=\rho-\frac12-\frac1{4D},
 \qquad M:=-C_0,
 \qquad K_0:=\frac1{4D}-D.
\end{equation}
For a fixed comparison assignment $O$, proving the theorem is
equivalent, by \eqref{eq:objective-pair-form}, to proving
\begin{equation}\label{eq:support-functional}
 \mathcal H:=C(W)-\rho C(O)-DQ_J\leq0.
\end{equation}

We first collect the parameter facts used below.

\medskip
\noindent
\textbf{Parameter facts.}
For every $k\geq1$, direct substitution of $D=1/(2h)$ gives
\begin{align}
 2h\left(C_0+Dk+\frac Ak\right)
 &=-\left[h^2-\left(1+\frac1k\right)h
              -\left(k+\frac1k+\frac1\alpha\right)\right].
 \label{eq:canonical-envelope-calculation}
\end{align}
The positive root of the quadratic in brackets is $h_k$.  Because
$h=h_n=\min_kh_k$, it follows that
\begin{equation}\label{eq:canonical-envelope}
 C_0+Dk+\frac Ak\geq0\quad(k\geq1),
 \qquad
 M=Dn+\frac An=\min_{k\geq1}\left(Dk+\frac Ak\right).
\end{equation}
For $\alpha\leq\tfrac12$, $h_2<h_1$, so $n\geq2$.  Indeed,
with $t=1/\alpha\geq2$,
\[
 h_1=1+\sqrt{t+3},\qquad
 h_2=\frac34+\sqrt{t+\frac{49}{16}},
\]
and the difference of the two square roots is less than $1/4$.
Moreover, $h_k\geq3$ for every $k$: evaluating the defining
quadratic at $3$ gives
\[
 3^2-3\left(1+\frac1k\right)
   -\left(k+\frac1k+\frac1\alpha\right)
 \leq4-k-\frac4k\leq0.
\]
Thus
\begin{equation}\label{eq:D-basic-range}
 0<D\leq\frac16.
\end{equation}

The condition that $n$ minimizes $Dk+A/k$, together with
\eqref{eq:D-basic-range}, yields
\begin{equation}\label{eq:active-D-range}
 \frac1{10}\leq D\leq\frac16\quad(n=2),
\end{equation}
and, for $n\geq3$,
\begin{equation}\label{eq:active-D-range-large-n}
 \frac1{2(n(n+1)-1)}\leq D
 \leq\frac1{2(n(n-1)-1)}.
\end{equation}
For example, $Dn+A/n\leq D(n+1)+A/(n+1)$ is equivalent to
$D\geq1/[2(n(n+1)-1)]$, and comparison with $n-1$ gives the
upper endpoint.  For $n=2$, comparison with $3$ gives $D\geq1/10$,
while \eqref{eq:D-basic-range} gives the stated upper endpoint.

The equality in \eqref{eq:canonical-envelope} also gives
\begin{equation}\label{eq:a-explicit}
 a=\frac{n-1}{4n}+\frac1{8D}
       -\frac D2\left(n+\frac1n\right).
\end{equation}
Set
\begin{equation}\label{eq:R-def}
 R:=\sqrt{\frac aD},\qquad \tau:=\frac1R=\sqrt{\frac Da}.
\end{equation}
We shall use
\begin{equation}\label{eq:R-sandwich}
 n\leq R\leq\frac{an}{A}.
\end{equation}
Here are full algebraic checks.  From \eqref{eq:a-explicit},
\begin{align}
 8Dn(a-Dn^2)
 &=n+2D(n-1)-4D^2(2n^3+n^2+1),\label{eq:R-lower-poly}\\
 8(aDn^2-A^2)
 &=n^2-2+2D(n(n-1)-4)-4D^2(n^3+n+2).
 \label{eq:R-upper-poly}
\end{align}
For $n=2$, the right sides are respectively $2+2D-84D^2$
and $2-4D-48D^2$.  Both are decreasing on
$[1/10,1/6]$ and both vanish at \(D=1/6\), so they are nonnegative
throughout that interval.  For $n\geq3$, put $B=n^2-n-1$.  Since
$D\leq1/(2B)$, the right side of \eqref{eq:R-lower-poly} is at least
\[
 n-\frac{2n^3+n^2+1}{B^2}\geq0,
\]
because
\[
 nB^2-(2n^3+n^2+1)
 =n^5-2n^4-3n^3+n^2+n-1
 =n^3(n-3)(n+1)+n^2+n-1>0.
\]
Likewise, the right side of \eqref{eq:R-upper-poly} is at least
\[
 n^2-2-\frac{n^3+n+2}{B^2}\geq0;
\]
after multiplication by $B^2$, the last inequality is
\[
 n^6-2n^5-3n^4+5n^3+3n^2-5n-4>0,
\]
which is immediate from
$n^4(n-3)(n+1)\geq0$ and
$5n^3+3n^2-5n-4>0$.  The first inequality in
\eqref{eq:R-sandwich} is equivalent to
$a-Dn^2\geq0$; the second, after squaring, is equivalent to
$aDn^2-A^2\geq0$.  This proves \eqref{eq:R-sandwich}.

We shall also use
\begin{equation}\label{eq:remaining-parameter-inequalities}
 M\leq A,\qquad
 \frac1n+\frac{D(n-1)}A\leq\frac34,\qquad
 4Da<1.
\end{equation}
For $n=2$, the first two inequalities follow directly from
$D\leq1/6$.  For $n\geq3$, the upper active-range inequality is
equivalent to $D\leq A/[n(n-1)]$.  It gives
\[
 \frac1n+\frac{D(n-1)}A\leq\frac2n\leq\frac23
\]
and
\[
 Dn\leq\frac A{n-1}\leq\frac{A(n-1)}n,
\]
the latter being $M=Dn+A/n\leq A$.  Finally,
\eqref{eq:a-explicit} gives
\[
 4Da=\frac12+\frac{D(n-1)}n
        -2D^2\left(n+\frac1n\right)
 <\frac12+D\leq\frac23.
\]

\medskip
\noindent
\textbf{1. Threshold decomposition and cell merging.}
Let $T$ be the minimum final machine load in $W$.  If $T=0$,
then, after zero-size jobs are discarded, there are fewer jobs than
machines and every job is alone in $W$; this assignment is optimal.
Thus assume $T>0$ and scale all processing times so that $T=1$.

On every $W$-machine whose final load is greater than $1$, designate
the last job assigned to that machine as its \emph{tail}.  If its size is
$p$ and the preceding load on that machine is $s$, then $s\leq1$:
when the tail was assigned, that machine had minimum current load.  The
machine whose final load is $1$ had current load at most $1$ at that
moment, so the minimum current load, and hence $s$, was at most $1$.
Call the preceding jobs the prefix of that row.  A row of final load
exactly $1$ is called a base row and all its jobs are base jobs.

Represent prefix and base jobs by the intersections of their $W$-rows
and their $O$-machines.  If two such jobs of sizes $u,v$ occupy the
same intersection, replace them algebraically by a single job of size
$u+v$.  The collision terms $C(W)$ and $C(O)$ both decrease by
$uv$, while $Q_J$ increases by $2uv$.  Hence $\mathcal H$ changes by
\begin{equation}\label{eq:cell-merge-change}
 (-uv)-\rho(-uv)-2Duv
 = (\rho-1-2D)uv=2\beta uv\geq0.
\end{equation}
We are maximizing a relaxation of $\mathcal H$, so it is harmless that the
merged algebraic configuration need no longer be generated by the
original list.  After all such mergers, every prefix/base intersection
contains at most one cell.

\medskip
\noindent
\textbf{2. Flattening all prefixes.}
Fix an overloaded row, condition on every other variable, and write its
prefix cell vector as $sv$, where $v\geq0$ and
$\sum_gv_g=1$.  If the current prefix load is $s=0$, choose any
probability vector $v$; the original point is still represented, and
the ensuing maximization is a valid relaxation.  Let $\tau_0$ be the
$O$-machine containing the
tail, put $x=v_{\tau_0}$, let $z\geq0$ be the total size of all other
jobs in that $O$-machine, and let
$d=\sum_gv_gd_g\geq0$, where $d_g$ is the load of the other jobs in
$O$-machine $g$.  Up to terms independent of $s,p$, the support
functional equals
\begin{equation}\label{eq:Jsp}
 J(s,p)=-Dp^2+sp(1-\rho x)-\rho zp-\rho ds
 +s^2\left(\frac12-A\lVert v\rVert_2^2\right),
 \qquad 0\leq s\leq1,\qquad p\geq1-s.
\end{equation}
For fixed $s$, the unconstrained maximizer in $p$ is
\[
 p^*(s)=\frac{s(1-\rho x)-\rho z}{2D}.
\]
On the interval where $p^*(s)>1-s$, maximizing over $p$ produces a
quadratic in $s$ with leading coefficient
\begin{equation}\label{eq:K-leading}
 K=\frac12-A\lVert v\rVert_2^2
       +\frac{(1-\rho x)^2}{4D}.
\end{equation}
We now prove $K\geq0$.  Since
$\lVert v\rVert_2^2\leq x^2+(1-x)^2$, multiplication by $4D$
shows that $4DK$ is at least
\begin{equation}\label{eq:P-x}
 P(x):=(\rho^2-4D-8D^2)x^2
 +(4D+8D^2-2\rho)x+1-4D^2.
\end{equation}
Its leading coefficient is positive, because
\[
 \rho^2-4D-8D^2
 \geq(1+2D)^2-4D-8D^2=1-4D^2>0.
\]
Its discriminant is $16D E$, where
\begin{align*}
 E&=D\rho^2-\rho(1+2D)+(1-D)(1+2D)^2\\
  &=D\bigl(\rho-(1+2D)\bigr)
     \left(\rho-\frac{(1+2D)(1-D)}D\right).
\end{align*}
The first factor in the product is nonnegative because
$\rho-1=D/\alpha\geq2D$.  The second is nonpositive.  Indeed, that
claim is equivalent to
$D^2(1/\alpha+2)\leq1$, while
\[
 D^2\left(\frac1\alpha+2\right)
 =\frac{1/\alpha+2}{4h^2}\leq\frac14,
\]
because every $h_k$, and hence $h$, is larger than
$\sqrt{1/\alpha+2}$.  Thus the discriminant is nonpositive.
Since $P(0)=1-4D^2>0$, $P(x)\geq0$ and therefore $K\geq0$.

The set
\[
 I=\{s\in[0,1]:p^*(s)>1-s\}
\]
is an interval and does not contain $0$.  On $I$, the optimized
value is convex by $K\geq0$, so its maximum over the closure of $I$
is attained either at $s=1$ or where $p^*(s)=1-s$.  Outside $I$,
the constrained maximizer already has $p=1-s$.  Consequently a global
maximizer of \eqref{eq:Jsp} has either $s=1$ or $p=1-s$.  Such a
maximizer exists because the domain in $s$ is compact and the term
$-Dp^2$ is coercive in $p$.

Apply this replacement successively to the finitely many overloaded
rows, always conditioning on the current values of all other rows.  The
support functional never decreases.  If $s=1$, the prefix has become
flat.  If $p=1-s$, the whole row has load $1$; reclassify the former
tail as a base job and merge it with an existing cell in the same
intersection if necessary, using \eqref{eq:cell-merge-change}.  Rows
already treated remain flat when later rows are changed.  After this
finite induction, every prefix/base row has total load $1$.

\medskip
\noindent
\textbf{3. Removing tail gains from nonempty optimal columns.}
Let $X=(x_{ig})$ be the resulting matrix of prefix/base cells.  It has
$m$ rows, $m$ columns, and every row sum is $1$.  Put
\[
 L_g:=\sum_i x_{ig},\qquad Q:=\sum_{i,g}x_{ig}^2.
\]
If tails of sizes $t_1,\dots,t_s$ lie in $O$-column $g$, their
contribution to $\mathcal H$ is
\begin{equation}\label{eq:tail-column-contribution}
 \sum_i(t_i-Dt_i^2)-\rho L_g\sum_i t_i
 -\rho\sum_{i<j}t_it_j.
\end{equation}
Replacing two tails $u,v$ by one algebraic tail of size $u+v$
increases \eqref{eq:tail-column-contribution} by
$(\rho-2D)uv>0$.  Thus we may retain at most one potential tail in
each $O$-column.  Optimizing its nonnegative size gives
\begin{equation}\label{eq:gL}
 g(L):=\max_{t\geq0}\{t-Dt^2-\rho Lt\}
      =\frac{(1-\rho L)_+^2}{4D}.
\end{equation}
Allowing a potential tail independently in every column is an upper
relaxation.  The prefix/base contribution is obtained from
$C(W)=m/2-Q/2$ and
$C(O)=\frac12\sum_gL_g^2-Q/2$.  Hence
\begin{equation}\label{eq:H-flat-upper}
 \mathcal H\leq \Phi(X):=\frac m2+\beta Q-a\sum_gL_g^2+
                    \sum_g g(L_g).
\end{equation}

The row-stochastic matrix polytope is compact and $\Phi$ is continuous;
choose a global maximizer $X$.  We claim that no column load lies
strictly between $0$ and $1/\rho$.  Suppose
$0<L_b<1/\rho$, set $v_i=x_{ib}/L_b$, and, for $g\neq b$, set
\[
 w_{ig}:=\frac{x_{ig}}{1-x_{ib}}.
\]
The denominator is positive because $x_{ib}\leq L_b<1$.  For
$0\leq z\leq1/\rho$, define a row-stochastic matrix by
\begin{equation}\label{eq:column-variation}
 x_{ib}(z)=zv_i,\qquad
 x_{ig}(z)=(1-zv_i)w_{ig}\quad(g\neq b).
\end{equation}
The original matrix occurs at $z=L_b$, a strict interior point of this
segment.

The coefficient of $z^2$ in $Q(X(z))$ is
\[
 U:=\sum_i v_i^2\left(1+\sum_{g\neq b}w_{ig}^2\right)\leq2,
\]
because $\sum_iv_i=1$ and every $w_i$ is a probability vector.
Writing $r_g=\sum_iv_iw_{ig}$, the coefficient of $z^2$ in
$\sum_gL_g(z)^2$ is
\[
 V:=1+\sum_{g\neq b}r_g^2\leq2,
\]
because $\sum_{g\neq b}r_g=1$.  Thus the quadratic part of
$\beta Q-a\sum_gL_g^2$ has coefficient at least
$-2a=-\rho$.  On the whole interval,
$g(L_b(z))=(1-\rho z)^2/(4D)$, whose quadratic coefficient is
$\rho^2/(4D)$.  Every other $g(L_g(z))$ is convex, since $g$ is
convex and $L_g(z)$ is affine.  Therefore $\Phi(X(z))$ is strictly
convex: its quadratic coefficient is at least
\[
 \frac{\rho^2}{4D}-\rho>0,
\]
where the last inequality follows from $4D\leq2/3<\rho$.
An interior point of a nonconstant strictly convex function cannot be a
maximum on a segment, contradicting the choice of $X$.

It follows that every nonempty column has load at least $1/\rho$, so
its tail gain is zero.  Let $e$ be the number of empty columns.  Only
these columns contribute to $\sum_g g(L_g)$, each by $1/(4D)$.
Delete the empty columns and transpose the remaining matrix.  In the
new orientation there are
\begin{equation}\label{eq:c-r-d}
 c:=m\text{ unit-sum columns},\qquad
 r:=m-e\text{ rows},\qquad d:=e=c-r.
\end{equation}
Writing the row sums again as $L_i$, it remains, by
\eqref{eq:H-flat-upper}, to prove
\begin{equation}\label{eq:transport-target}
 a\sum_{i=1}^rL_i^2-\beta Q-\frac c2-\frac d{4D}\geq0.
\end{equation}

\medskip
\noindent
\textbf{4. Purification and integer peeling.}
For a nonnegative matrix with $c$ unit-sum columns and prescribed row
sums $L=(L_i)$, let
\[
 V_c(L):=\max\sum_{i,j}x_{ij}^2.
\]
The feasible transportation polytope is compact, so this maximum is
attained whenever the margins are feasible.
We need the following exact identity.

\smallskip
\noindent
\emph{Purification lemma.}
If $L_i\geq1$, one can make some column pure in row $i$, preserving
all margins and without decreasing the square sum.  Consequently, if
$z_i=\lfloor L_i\rfloor$, $f_i=L_i-z_i$, and
$N=\sum_i z_i$, then
\begin{equation}\label{eq:purification-identity}
 V_c(L)=N+V_{c-N}(f).
\end{equation}

To prove the lemma, choose a target column $b$ maximizing
$a_0=x_{ib}$, and put $R_0=1-a_0$.  If $R_0=0$, the column is
already pure.  Otherwise choose numbers
$0\leq d_j\leq x_{ij}$, $j\neq b$, with
$\sum_{j\neq b}d_j=R_0$.  This is possible because
$L_i-a_0\geq1-a_0=R_0$.  If $y_h=x_{hb}$ for $h\neq i$, choose a
nonnegative transport $\varepsilon_{hj}$ with row sums $y_h$ and
column sums $d_j$; for example,
$\varepsilon_{hj}=y_hd_j/R_0$.  Make column $b$ pure in row $i$,
remove $d_j$ from cell $(i,j)$, and add
$\varepsilon_{hj}$ to cell $(h,j)$.  All margins are preserved.
Since $x_{ij}\leq a_0$ for every donor column, the change in square sum
is at least
\begin{align*}
 &(1-a_0^2)-\sum_hy_h^2
 +\sum_j(d_j^2-2a_0d_j)+\sum_{h,j}\varepsilon_{hj}^2\\
 &\qquad=R_0^2-\sum_hy_h^2+\sum_jd_j^2
              +\sum_{h,j}\varepsilon_{hj}^2\geq0,
\end{align*}
because $\sum_hy_h=R_0$.  Repeating this operation peels exactly
$z_i$ pure columns from row $i$, for every $i$.  Starting with a
maximizer proves the upper inequality in
\eqref{eq:purification-identity}; adjoining the $N$ pure columns to a
maximizing residual matrix proves the reverse inequality.

\medskip
\noindent
\textbf{5. The scalar partition lemma and the basic residual bound.}
Define
\begin{equation}\label{eq:phi-def}
 \phi(x):=
 \begin{cases}
 2\sqrt{aD}\,x-\beta x^2,&0\leq x\leq\tau,\\
 Ax^2+D,&\tau\leq x\leq1.
 \end{cases}
\end{equation}
The two expressions and their first derivatives agree at $\tau$.

\smallskip
\noindent
\emph{Scalar partition lemma.}
For every finite partition $x_1+\cdots+x_s=1$, $x_j\geq0$,
\begin{equation}\label{eq:scalar-partition}
 \sum_j\phi(x_j)\geq M.
\end{equation}

Discard zero parts.  Among parts strictly below $\tau$, hold the sum of
two parts fixed.  Their total contribution is a concave quadratic in one
of the two parts, so moving to an endpoint makes one part zero or
$\tau$ and does not increase the sum.  Repetition leaves at most one
part $y\in[0,\tau]$ below the threshold.  If there are $k$ remaining
parts, all at least $\tau$, convexity of $Ax^2+D$ allows us to replace
them by $k$ equal parts $(1-y)/k$.  Thus it suffices to consider
\begin{equation}\label{eq:Fky}
 F_k(y)=Dk+\frac Ak(1-y)^2+2\sqrt{aD}\,y-\beta y^2,
 \qquad
 0\leq y\leq u_k:=\min\{\tau,1-k\tau\}.
\end{equation}
Only integers with $k\leq R$ occur.  Expanding,
\begin{equation}\label{eq:Fky-expanded}
 F_k(y)=Dk+\frac Ak
 +2\left(\sqrt{aD}-\frac Ak\right)y
 +\left(\frac Ak-\beta\right)y^2.
\end{equation}
Its minimum on the feasible interval is attained at an endpoint.  To
see this, the assertion is immediate if the quadratic coefficient is
nonpositive.  If it is positive and the derivative at zero is
nonnegative, the function is increasing.  In the remaining case its
vertex is
\[
 y_*=\frac{A/k-\sqrt{aD}}{A/k-\beta}.
\]
The negative derivative condition says
$k<AR/a$.  By \eqref{eq:R-sandwich}, $AR/a\leq n$, so
$k\leq n-1\leq R-1$.  Moreover
\[
 y_*\geq\tau
 \quad\Longleftrightarrow\quad
 \frac A k(1-\tau)\geq A\tau
 \quad\Longleftrightarrow\quad k\leq R-1.
\]
Thus the vertex lies to the right of the feasible interval and again the
minimum is at an endpoint.

At $y=0$, \eqref{eq:canonical-envelope} gives
$F_k(0)=Dk+A/k\geq M$.  If $u_k=\tau$, the endpoint partition can be
viewed as $k+1$ parts on the convex branch (with $\tau$ belonging to
either branch), and Jensen's inequality gives
\[
 F_k(\tau)\geq D(k+1)+\frac{A}{k+1}\geq M.
\]
If $u_k=1-k\tau<\tau$, all endpoint parts are at most $\tau$.  For
$0<x\leq\tau$,
\begin{align}
 \frac{\phi(x)}x
 &\geq2\sqrt{aD}-\beta\tau
   =DR+\frac AR
   \geq Dn+\frac An=M. \label{eq:phi-linear-lower}
\end{align}
The last inequality follows because
\[
 DR+\frac AR-\left(Dn+\frac An\right)
 =(R-n)\left(D-\frac{A}{nR}\right)\geq0;
\]
the two factors are nonnegative by \eqref{eq:R-sandwich} and
$a=DR^2$.  Hence the sum over an all-small endpoint partition is at
least $M\sum_jx_j=M$.  This proves
\eqref{eq:scalar-partition}.

We now deduce a matrix inequality.  Let a matrix have $q$ unit-sum
columns, $r$ rows with arbitrary row sums $s_i\geq0$, and square sum
$Q$.  Then
\begin{equation}\label{eq:basic-residual}
 a\sum_i s_i^2-\beta Q+Dr+C_0q\geq0.
\end{equation}
For each nonzero row, distribute $as_i^2+D$ among its cells
proportionally to cell size.  A cell $x\leq s_i$ receives, after its
$-\beta x^2$ term is included, at least
\[
 x\inf_{s\geq x}\left(as+\frac Ds\right)-\beta x^2
 =\phi(x).
\]
Indeed, the unconstrained minimum occurs at $s=\tau$, giving the first
branch in \eqref{eq:phi-def}; when $x\geq\tau$, it occurs at $s=x$,
giving $ax^2+D-\beta x^2=Ax^2+D$.  Applying the scalar partition lemma
in every column gives a total of at least $qM$.  Since $C_0=-M$, this
cancels $C_0q$.  Every zero row retains its unused positive $D$, so
\eqref{eq:basic-residual} follows.

\medskip
\noindent
\textbf{6. A strengthened residual bound for fractional row sums.}
Suppose now that every row sum satisfies $0\leq s_i<1$.  We prove
\begin{equation}\label{eq:strengthened-residual}
 E:=a\sum_i s_i^2-\beta Q+Dr+C_0q
 \geq\sum_i(\rho s_i-K_0)_+.
\end{equation}
Put
\[
 G(s):=as^2+D-(\rho s-K_0)_+,
 \qquad H_0:=\frac1{4D}-a.
\]
For a cell of size $x$ in a row of sum $s<1$, proportional
distribution of $G(s)$, followed by inclusion of $-\beta x^2$, gives,
for $0<x<1$, the lower cell charge
\begin{equation}\label{eq:psi-def-inf}
 \psi(x):=x\inf_{x\leq s<1}\frac{G(s)}s-\beta x^2.
\end{equation}
Set $\psi(0):=0$ and define \(\psi(1)\) by continuous extension.

Let $x_0$ be the larger root of
\begin{equation}\label{eq:x0-equation}
 ax^2-H_0x+D=0.
\end{equation}
We first establish
\begin{equation}\label{eq:x0-range}
 \frac{\sqrt3}{2}\leq x_0\leq1.
\end{equation}
Let $s_0=\sqrt3/2$, $C=s_0^2+s_0$, and define
\[
 u:=\frac{s_0(1-s_0)}8,
 \quad v_n:=\frac C2\left(n+\frac1n\right)-1,
 \quad w_n:=\frac{C(n-1)}{4n}.
\]
Substitution of \eqref{eq:a-explicit} into the left side of
\eqref{eq:x0-equation} at $s_0$ gives exactly
\begin{equation}\label{eq:x0-s0-eval}
 w_n-\frac uD-v_nD.
\end{equation}
Here $u,v_n>0$, and $w_n^2\leq4uv_n$.  For completeness, for
$n=2,3,4$, the respective values of $4uv_n-w_n^2$ are
\[
 \frac{483-268\sqrt3}{1024},\qquad
 \frac{95-52\sqrt3}{192},\qquad
 \frac{2163-1132\sqrt3}{4096},
\]
all positive.  For $n\geq5$, $w_n\leq C/4$ and
$v_n\geq(13/5)C-1$, whence
\[
 4uv_n-w_n^2
 \geq4u\left(\frac{13}{5}C-1\right)-\frac{C^2}{16}
 =\frac{687-380\sqrt3}{1280}>0.
\]
The arithmetic--geometric mean inequality now makes
\eqref{eq:x0-s0-eval} nonpositive.  Thus the quadratic in
\eqref{eq:x0-equation} has real roots and $s_0$ lies between them.

At $x=1$, the same quadratic equals
\[
 a-H_0+D=A-M\geq0,
\]
where the identity follows from \eqref{eq:C0-M-K0} and the inequality
from \eqref{eq:remaining-parameter-inequalities}.  Also
$H_0>0$ by $4Da<1$, so both roots are positive.  Since
$s_0<1$, the larger root lies between $s_0$ and $1$, proving
\eqref{eq:x0-range}.

We next compute $\psi$.  The penalty in $G$ starts at
$\sigma:=K_0/\rho$.  The identity $K_0-\rho=M-A\leq0$ gives
$0<\sigma\leq1$.  On the unpenalized branch,
\[
 \frac{G(s)}s=as+\frac Ds;
\]
on the penalized branch,
\[
 \frac{G(s)}s=as-\rho+\frac1{4Ds}.
\]
The latter is decreasing up to $s=1$, because $4Da<1$, and its
infimum is $H_0$.  Equation \eqref{eq:x0-equation} gives
\begin{equation}\label{eq:sigma-x0}
 \sigma-x_0=\frac{1-x_0}{2}
       \left(1+\frac{D}{ax_0}\right)\geq0.
\end{equation}
Moreover, the discriminant of \eqref{eq:x0-equation} is nonnegative, so
\(H_0\geq2\sqrt{aD}\).  If \(x_-\) is its smaller root, Vieta's
formula gives $x_-x_0=D/a=\tau^2$.  Since
$x_0\geq\sqrt3/2>1/2\geq\tau$, it follows that
\[
 x_-\leq\tau\leq x_0.
\]
For \(x\leq\tau\), the infimum of the unpenalized branch is
\(2\sqrt{aD}\leq H_0\).  For \(\tau\leq x\leq x_0\), its infimum is
\(ax+D/x\), and
\[
 ax+\frac Dx\leq H_0
 \quad\Longleftrightarrow\quad ax^2-H_0x+D\leq0.
\]
For \(x_0\leq x<1\), this inequality is reversed, while the decreasing
penalized branch has infimum \(H_0\).  Together with
\(x_0\leq\sigma\leq1\), these comparisons give
\begin{equation}\label{eq:psi-piecewise}
 \psi(x)=
 \begin{cases}
 \phi(x),&0\leq x\leq x_0,\\
 H_0x-\beta x^2,&x_0\leq x\leq1.
 \end{cases}
\end{equation}
When $\sigma=1$, the same conclusion follows by continuity; in that
case $x_0=1$.

By \eqref{eq:x0-range} and
\eqref{eq:remaining-parameter-inequalities},
\begin{align}
 \psi(x_0)=Ax_0^2+D
 &\geq \frac{3A}{4}+D
 \geq \frac An+Dn=M.\label{eq:psi-x0-lower}
\end{align}
The second inequality is precisely
$1/n+D(n-1)/A\leq3/4$.  Also
\begin{equation}\label{eq:psi-one}
 \psi(1)=H_0-\beta=M+D.
\end{equation}
The second branch in \eqref{eq:psi-piecewise} is concave, so
$\psi(x)\geq M$ throughout $[x_0,1]$.  In addition,
$\psi(x)\geq0$ on all of $[0,1]$: this is clear on the second branch,
and on the first it follows from \eqref{eq:phi-def} and
$2\sqrt{aD}-\beta\tau=(a+A)\tau>0$.

Consider one unit-sum column.  If all its cells are at most $x_0$,
then \eqref{eq:psi-piecewise} and the scalar partition lemma give total
charge at least $M$.  Otherwise, since $x_0\geq\sqrt3/2>1/2$, there
is a unique cell larger than $x_0$; that cell alone contributes at
least $M$, and all other cells contribute nonnegatively.  Thus every
column has total $\psi$-charge at least $M$.  Summing over columns,
using $C_0=-M$, and retaining the unused $D$ from every zero row proves
\eqref{eq:strengthened-residual}.

\medskip
\noindent
\textbf{7. Completion of the transport inequality.}
Return to \eqref{eq:transport-target}.  Write
\[
 L_i=z_i+f_i,\qquad z_i=\lfloor L_i\rfloor,\qquad0\leq f_i<1.
\]
Because $\sum_iL_i=c$, the number
\[
 q:=\sum_i f_i=c-\sum_i z_i
\]
is an integer.  By the purification identity,
\begin{equation}\label{eq:Q-purified-upper}
 Q\leq V_c(L)=c-q+V_q(f).
\end{equation}
Since $\beta\geq0$, it is enough to use the right side in place of
$Q$.  For fixed $f$, choose a nonnegative integer vector $z$ with
$\sum_i z_i=c-q$ minimizing
$\sum_i(z_i+f_i)^2$.  This can only lower the left side of
\eqref{eq:transport-target}.  Such a minimizer has one of the following
two forms:
\begin{enumerate}
 \item every $z_i\geq1$; or
 \item every $z_i\in\{0,1\}$.
\end{enumerate}
Indeed, if $z_i=0$ and $z_j\geq2$, moving one unit from $j$ to
$i$ changes the square sum by
\[
 2(1-z_j+f_i-f_j)<0,
\]
contradicting minimality.

In the first case, put $s_i=L_i-1=z_i-1+f_i$.  There are
$r=c-d$ rows and $\sum_i s_i=d$.  Applying the purification identity
to $s$ shows
\[
 c-q+V_q(f)=r+V_d(s).
\]
Let a $d$-column residual matrix attain $V_d(s)$.  Substitution in
the left side of \eqref{eq:transport-target} gives
\begin{align*}
 &a\sum_i(1+s_i)^2-\beta(r+V_d(s))
       -\frac c2-\frac d{4D}\\
 &\qquad=a\sum_i s_i^2-\beta V_d(s)+Dr+C_0d,
\end{align*}
where we used $\sum_i s_i=d$, $c=r+d$, $2a=\rho$, and the
definitions of $A,C_0$.  This is nonnegative by the basic residual
bound \eqref{eq:basic-residual}.

In the second case, let
\[
 Z:=\{i:z_i=0\},\qquad t:=|Z|,\qquad
 F_0:=\sum_{i\in Z}f_i.
\]
The number of one entries in $z$ is $c-q$, so
\begin{equation}\label{eq:t-q-d}
 t=r-(c-q)=q-d.
\end{equation}
Let a $q$-column fractional matrix with row sums $f$ attain
$V_q(f)$, and denote by
\[
 E:=a\sum_i f_i^2-\beta V_q(f)+Dr+C_0q
\]
the left side of its strengthened residual inequality.  A direct
expansion, using \eqref{eq:t-q-d}, gives
\begin{equation}\label{eq:final-expansion}
 a\sum_i(z_i+f_i)^2-\beta(c-q+V_q(f))
       -\frac c2-\frac d{4D}
 =E+tK_0-\rho F_0.
\end{equation}
By \eqref{eq:strengthened-residual},
\[
 E\geq\sum_i(\rho f_i-K_0)_+
 \geq\sum_{i\in Z}(\rho f_i-K_0)
 =\rho F_0-tK_0.
\]
Together with \eqref{eq:final-expansion}, this proves
\eqref{eq:transport-target} in the second case as well.

Thus \eqref{eq:transport-target} always holds.  Reversing the reductions
in Sections 3, 2, and 1 gives $\mathcal H\leq0$ in
\eqref{eq:support-functional}.  Therefore
$\gamma_{\alpha}(W)\leq\rho \gamma_{\alpha}(O)$.  Taking $O$ optimal completes the
proof.
\end{proof}

\section{Matching instances and the exact ratio}

\begin{proposition}[Matching finite instances]
\label{prop:matching-instance}
For every $0<\alpha\leq1/2$ and every index $k$ attaining
$H(\alpha)=h_k(\alpha)$, there is an instance on $k+1$ machines and
a list schedule whose ratio is exactly
$\rho_k(\alpha)=1+1/(2\alpha h_k(\alpha))$.
\end{proposition}

\begin{proof}
Use $m=k+1$ machines, and take $k$ jobs of size
$1$, $k$ jobs of size $1/k$, and one job of size $h_k$.  A list can
first put the $k$ unit jobs and one small job on the empty machines,
then put the remaining small jobs on the small-job machine, and finally
schedule the long job.  Immediately before the long job, every machine
has load \(1\), so the final minimum-load tie is harmless: whichever
machine is selected, the long job's collision contribution is \(h_k\).
The list schedule's value is
\[
 F_W=\alpha\left(h_k^2+k+\frac1k\right)
       +\frac{k-1}{2k}+h_k.
\]
Pairing each unit job with one small job and putting the long job alone
has value
\[
 F_O=\alpha\left(h_k^2+k+\frac1k\right)+1.
\]
Since $h_k$ satisfies
\[
 h_k^2-\frac{k+1}{k}h_k
       -\left(k+\frac1k+\frac1\alpha\right)=0,
\]
one obtains the ratio explicitly as follows.  With
$b_k=(k+1)/(2k)$,
\[
 F_W-F_O=h_k-b_k,
 \qquad
 F_O=2\alpha h_k(h_k-b_k),
\]
where the second identity follows by substituting the displayed
quadratic into $F_O$.  Therefore
$F_W/F_O=1+1/(2\alpha h_k)=\rho_k$.
The displayed comparison schedule is optimal.  By
\eqref{eq:objective-load-form}, $Q_J$ is fixed by the instance, so it
suffices to minimize the sum of squared machine loads.  First,
\[
 h_k>\sqrt{k+\frac1k+\frac1\alpha}
 \geq\sqrt{k+\frac1k+2}
 =\sqrt{k}+\frac1{\sqrt{k}}\geq1+\frac1k.
\]
If total load \(x>0\) is colocated with the long job, choose one attached
job of size \(p\).  The least-loaded one of the other \(k\) machines has
load at most \((k+1-x)/k\leq1+1/k<h_k\).  Moving \(p\) to that machine
strictly decreases the sum of squared loads: if the destination load is
$L$, the change is
\[
 2p\bigl[L-(h_k+x-p)\bigr]<0,
\]
because the long machine still has load $h_k+x-p\geq h_k>L$ after
$p$ is removed.  Hence the long job is alone in every optimum.  The
other $k$ machines then have total load $k+1$, and the displayed
assignment attains the square-sum lower bound by making all their loads
equal to $1+1/k$.
\end{proof}

\begin{proposition}[The active integer]
\label{prop:active-index}
For $k\geq2$, define
\[
 T_k:=k^4+2k^3-2k^2-5k+1.
\]
For each integer $k\geq2$, the index $k$ minimizes $h_\ell(\alpha)$
over all $\ell\geq1$ if and only if
\[
 \begin{cases}
 \displaystyle \frac1{15}\leq\alpha\leq\frac12,
   & k=2,\\[6pt]
 \displaystyle \frac1{T_k}\leq\alpha\leq\frac1{T_{k-1}},
   & k\geq3.
 \end{cases}
\]
At a shared endpoint the two adjacent indices are both active.
These intervals cover all $0<\alpha\leq1/2$.
\end{proposition}

\begin{proof}
For an integer $\ell\geq1$, introduce
\[
 q_\ell(x):=
 x^2-\frac{\ell+1}{\ell}x
 -\left(\ell+\frac1\ell+\frac1\alpha\right).
\]
Its positive root is $h_\ell$.  Since $q_k(h_k)=0$, subtraction gives
\begin{equation}
 q_\ell(h_k)
 =\frac{(\ell-k)(h_k+1-k\ell)}{k\ell}.
 \label{eq:root-comparison}
\end{equation}
The other root of $q_\ell$ is negative.  Therefore, for $h_k>0$,
\[
 h_k\leq h_\ell
 \quad\Longleftrightarrow\quad
 q_\ell(h_k)\leq0.
\]
For $\ell>k$, equation \eqref{eq:root-comparison} requires
$h_k\leq k\ell-1$, and the strongest such requirement is obtained at
$\ell=k+1$.  For $\ell<k$, it requires
$h_k\geq k\ell-1$, and the strongest requirement is obtained at
$\ell=k-1$.  Hence
\begin{equation}
 h_k\leq h_\ell\quad\text{for every $\ell\geq1$}
 \quad\Longleftrightarrow\quad
 k(k-1)-1\leq h_k\leq k(k+1)-1.
 \label{eq:active-window}
\end{equation}
The proof of Theorem~\ref{thm:exact-wlpt-alpha} already showed
$h_2<h_1$, so $k=1$ is never active.

The transition between $k$ and $k+1$ occurs when
$h_k=k(k+1)-1$.  Substitution in the defining quadratic
\[
 h_k^2-\frac{k+1}{k}h_k
 =k+\frac1k+\frac1\alpha
\]
gives
\[
 \frac1\alpha
 =k^4+2k^3-2k^2-5k+1=T_k.
\]
Moreover,
\[
 T_{k+1}-T_k=4k^3+12k^2+6k-4>0
 \qquad(k\geq2),
\]
and $T_2=15$, while $T_k\to\infty$.  For fixed $k$, the root $h_k$
is strictly increasing in $1/\alpha$.  Thus
\eqref{eq:active-window} says, for $k\geq3$, that $k$ is active
exactly when
\[
 T_{k-1}\leq\frac1\alpha\leq T_k,
\]
which is the stated interval after taking reciprocals.  For $k=2$, the
lower side of the active window is automatic on the present domain and
the upper side gives $1/\alpha\leq T_2=15$.  The monotonicity and
$T_k\to\infty$ also prove that the intervals cover
$(0,1/2]$ and meet only at the stated adjacent transitions.
\end{proof}

\begin{corollary}[Exact worst-list ratio]
\label{cor:exact-ratio}
For every $0<\alpha\leq1/2$,
\[
\boxed{\;
\sup_I\frac{\WLPT_\alpha(I)}{\OPT_\alpha(I)}
 =\max_{k\geq2}\left(1+\frac1{2\alpha h_k(\alpha)}\right)
 =1+\frac1{2\alpha H(\alpha)}.
\;}
\]
Here the supremum ranges over every finite positive-processing-time
instance and every finite number of identical machines.
In fact, under the tie convention stated in Section~1, a finite
matching instance attains the displayed value.
\end{corollary}

\begin{proof}
Theorem~\ref{thm:exact-wlpt-alpha} proves the upper bound for every
instance and every list schedule.  Select an active integer from
Proposition~\ref{prop:active-index}; Proposition
\ref{prop:matching-instance} supplies a finite instance attaining the
same value.  The two inequalities coincide.
\end{proof}

\begin{remark}[Why a fixed $k$ is not a universal formula]
An individual expression $\rho_k(\alpha)$ is a universal upper bound
only when $k$ is active at that value of $\alpha$.  Outside its active
interval it is merely the value of the corresponding lower-bound
construction and is smaller than the true worst-case ratio.  The
envelope $1+1/(2\alpha H(\alpha))$, or equivalently the piecewise
choice in Proposition~\ref{prop:active-index}, is the universal bound.
\end{remark}

\begin{corollary}[Comparison with the benchmark bound]
\label{cor:strict-improvement}
For $0<\alpha<1/2$,
\[
 1+\frac1{2\alpha H(\alpha)}
 <1+\frac1{6\alpha}.
\]
At $\alpha=1/2$, both sides equal $4/3$, and no strict improvement is
possible.
\end{corollary}

\begin{proof}
Because $h_k$ is the positive root of $q_k$, evaluation at $x=3$
shows
\[
 h_k\geq3
 \quad\Longleftrightarrow\quad
 k+\frac4k+\frac1\alpha\geq6.
\]
For $k\geq2$, the arithmetic--geometric mean inequality gives
$k+4/k\geq4$, with equality only at $k=2$.  Also
$1/\alpha\geq2$, with equality only at $\alpha=1/2$.
Thus $H(\alpha)>3$ throughout $0<\alpha<1/2$, proving the strict
inequality.  At $\alpha=1/2$, $h_2=3$, so the exact ratio is $4/3$.
Proposition~\ref{prop:matching-instance} attains it, ruling out a
smaller universal upper bound.
\end{proof}

\begin{remark}[Endpoint check and provenance]
At $\alpha=1/2$, identity \eqref{eq:objective-load-form} becomes
$F_{1/2}(S)=\tfrac12\sum_iL_i(S)^2$.  The value $4/3$ is therefore
consistent with the sharp squared-load result of
Avidor, Azar, and Sgall~\cite{AvidorAzarSgall2001}.  That citation is
used only as an endpoint check.  To our knowledge, the parameterized
exact envelope is unpublished; it is presented here as a new,
self-contained derivation and is not attributed to that paper.
\end{remark}

\begin{thebibliography}{1}

\bibitem{AvidorAzarSgall2001}
A.~Avidor, Y.~Azar, and J.~Sgall,
\newblock Ancient and new algorithms for load balancing in the
  $\ell_p$ norm,
\newblock \emph{Algorithmica} 29(3):422--441, 2001.
\newblock \href{https://doi.org/10.1007/s004530010051}
  {doi:10.1007/s004530010051}.

\end{thebibliography}

\end{document}
